\par\Needspace{6\baselineskip}\medskip
\SourceLeaf{224}\begin{center}\hypertarget{native-224-article161}{}\pdfbookmark[2]{161.}{article-161}161.\end{center}
\emph{Casus primi} tres subdivisiones faciemus:

I. Si punctum $B$ cum $A$ vel cum puncto opposito coincidit, erit $\delta=0$
vel $=180^0$; $\gamma$, $\epsilon'$, $\epsilon''$ atque puncta $D'$, $D''$
indeterminata erunt; contra $\gamma'$, $\gamma''$, $\epsilon$ atque puncta $D$, $B^*$
determinata; punctum $C$ necessario coincidit cum $A$. Per ratiocinia, iis,
quae in art. 140 tradita sunt, analoga, facile elicietur aequatio haecce:
\[
0=n'\frac{\sin(z-\sigma)}{\sin z}\cdot\frac{R'\sin\delta'}{R''\sin\delta''}
\cdot\frac{\sin(A''D-\delta'')}{\sin(A'D-\delta'+\sigma)}-n''.
\]
Omnia itaque, quae in artt. 141, 142 exposita sunt, etiam huc transferre licebit,
si modo statuatur $a=0$, atque $b$ per ipsam aequationem 12 art. 140 determinetur,
quantitatesque $z$, $r'$, $\dfrac{n'r'}{n}\mathpunct{\raisebox{.45ex}{,}}$
$\dfrac{n'r'}{n''}$ perinde ut supra computabuntur. Iam simulac $z$ adeoque
situs puncti $C'$ innotuit, assignare licebit situm circuli maximi $CC'$, huius
intersectionem cum circulo maximo $A''B''$ i.\ e.\ punctum $C''$, et proin arcus
$CC'$, $CC''$, $C'C''$ sive $2f''$, $2f'$, $2f$: hinc denique habebitur
\[
r=\frac{n'r'}{n}\cdot\frac{\sin2f}{\sin2f'}\mathpunct{\raisebox{.45ex}{,}}\qquad
r''=\frac{n'r'}{n''}\cdot\frac{\sin2f''}{\sin2f'}\mathpunct{\cdot}
\]

II. Ad casum eum, ubi punctum $B''$ cum $A''$ vel cum puncto opposito coincidit,
omnia quae modo tradidimus transferre licet, si modo omnia, quae ad locum
primum spectant, cum iis, quae ad tertium referuntur, permutantur.

III. Paullo aliter vero casum eum tractare oportet, ubi $B'$ vel cum $A'$ vel
cum puncto opposito coincidit. Hic punctum $C'$ cum $A'$ coincidet; $\gamma'$,
$\epsilon$, $\epsilon''$ punctaque $D$, $D''$, $B^*$ indeterminata erunt: contra
assignari poterit intersectio circuli maximi $BB''$ cum ecliptica\footnote{Generalius, cum circulo maximo $AA''$: sed brevitatis caussa eum tantummodo casum hic consideramus, ubi ecliptica tamquam planum fundamentale accipitur.},
cuius longitudo ponatur $=l'+\pi$. Per ratiocinia, iis, quae in art. 140 evoluta
sunt, similia, eruetur aequatio
\[
0=n\frac{R\sin\delta\sin(A''D'-\delta'')}{R''\sin\delta''\sin(AD'-\delta)}
+n'r'\frac{\sin\pi}{R''\sin(l''-l'-\pi)}+n''.
\]
Designemus coëfficientem ipsius $n$, qui convenit cum $a$ art. 140, per eundem
\SourceLeaf{225}characterem $a$, coëfficientemque ipsius $n'r'$ per $\beta$: ipsum $a$
hic etiam per formulam
\[
a=-\frac{R\sin(l'+\pi-l)}{R''\sin(l''-l'-\pi)}
\]
determinare licet. Habemus itaque $0=an+\beta n'r'+n''$, qua aequatione cum his combinata
\[
P=\frac{n''}{n}\mathpunct{\raisebox{.45ex}{,}}\qquad
Q=2\left(\frac{n+n''}{n'}-1\right)r'^3,
\]
emergit
\[
\frac{\beta(P+1)}{P+a}r'^4+r'^3+\tfrac12Q=0,
\]
unde distantiam $r'$ elicere poterimus, siquidem non fuerit $\beta=0$, in quo
casu nihil aliud illinc sequeretur, nisi $P=-a$. Ceterum etiamsi non fuerit
$\beta=0$ (ubi ad casum tertium in art. sequ. considerandum delaberemur), tamen
semper $\beta$ quantitas perexigua erit, adeoque $P$ parum a $-a$ differre
debebit: hinc vero manifestum est, determinationem coëfficientis
$\dfrac{\beta(P+1)}{P+a}$ valde lubricam fieri, neque adeo $r'$ ulla praecisione
determinabilem esse.

Porro habebimus $\dfrac{n'r'}{n}=-\dfrac{P+a}{\beta}\mathpunct{\raisebox{.45ex}{,}}$
$\dfrac{n'r'}{n''}=-\dfrac{P+a}{\beta P}\mathpunct{:}$ dein simili modo ut in
art. 143 facile evolventur aequationes
\[
\begin{aligned}
r\sin\zeta&=\frac{n'r'}{n}\cdot\frac{\sin\gamma''}{\sin\epsilon'}\sin(l''-l')\\[.5em]
r''\sin\zeta''&=-\frac{n'r'}{n''}\cdot\frac{\sin\gamma}{\sin\epsilon'}\sin(l'-l)\\[.5em]
r\sin(\zeta-AD')&=r''P\frac{\sin\gamma''}{\sin\gamma}\sin(\zeta''-A''D'),
\end{aligned}
\]
e quarum combinatione cum aequatt. VIII et IX art. 143 quantitates $r$, $\zeta$,
$r''$, $\zeta''$ determinare licebit. Calculi operationes reliquae cum supra
descriptis convenient.
